% ============================================================
%  ch12 —— §8（原文 p.46–p.48）
% ============================================================
\routenote{p.46--p.48}

\section{附录：$\PP^1$ 上 $G$-丛模空间几何的命题 1.1 的证明}

本附录给出命题 1.1 的证明。

\begin{proof}
前两条断言的证明用与 [23, Proposition 3] 相同的论证。先假设基域 $k$ 代数闭。

\emph{A.} 设 $G=T$ 是环面。此时 (2) 已在 [24, Lemma 16] 中证明。下面证明 (1)。由于 $T$ 交换，对任何 $T$-丛 $P$ 有 $H^1(\PP^1,P\times^T\mathrm{Lie}(T))=H^1(\PP^1,\mathrm{Lie}(T))$。我们断言该群消没。

按构造，$\mathrm{Lie}(T)\subset\mathrm{Lie}([N^*](\mathbb{G}_m\times\PP^1))$ 是直和项。又 $\mathrm{Lie}([N]^*(\mathbb{G}_m\times\PP^1))=[N]^*\mathcal O_{\PP^1}$，而 $H^1(\PP^1,[N]^*(\mathcal O_{\PP^1}))=H^1(\PP^1,\mathcal O_{\PP^1})=0$。故对任何 $P$ 有 $H^1(\PP^1,P\times^T\mathrm{Lie}(T))=0$，于是 $T$-丛没有非平凡变形。由于 $k$ 代数闭，$H^1(k(\PP^1),T)=0$，即所有 $T$-丛都一般平凡。因此映射 $\oplus_{x\in\PP^1}T(K_x)/T(\mathcal O_x)\to\Bun_T(k)$ 是满的。按构造，同构 $\pi_0(\Bun_T)\cong\pi_1(T)^{\pi_1(\mathbb{G}_m)}$ 由 Kottwitz 同态 $\pi_0(\mathrm{Gr}_{T,x})\to X_*(T)^{\operatorname{Gal}(K^{\mathrm{sep}}_x/K_x)}$ 诱导。故对任何 $x$，映射 $T(K_x)\to\Bun_T(k)$ 是满的。这证明了环面情形的 (1)。最后，(1) 与 (2) 蕴涵环面情形的 (4)。

\emph{B.} 对半单单连通群 $G$，陈述 (1) 与 (2) 已在 [24, Theorem 4] 中证明（同样假设 $k$ 代数闭）。

\emph{C.} 设 $G|_{\PP^1\setminus\{0,\infty\}}$ 的导群 $G^{\mathrm{der}}|_{\PP^1\setminus\{0,\infty\}}$ 单连通，并记 $D|_{\PP^1\setminus\{0,\infty\}}:=G/G^{\mathrm{der}}|_{\PP^1\setminus\{0,\infty\}}$。Haines 与 Rapoport 证明 $\PP^1$ 上有群概形的诱导正合列
\[
1\to G^{\mathrm{der}}\to G\to D\to 1.
\]
由于任何 $D$-丛在 $\PP^1\setminus\{x\}$ 上平凡，且对 $G^{\mathrm{der}}$ 的所有形式也有同样的结论，这蕴涵 $G$ 的 (1)。此外，由（上引文献）知 $G(k((s)))\to D(k((s)))$ 是满的。于是 $\pi_0(\Bun_G)\cong\pi_0(\Bun_D)$，这证明了 $G$ 的 (2)。

\emph{D.} 对一般群 $G$，可取 $z$-扩张
\[
1\to Z\to G'\to G\to 1,
\]
使 $Z$ 是诱导环面且 $G'^{\mathrm{der}}$ 单连通。由 Tsen 定理 $H^2(\PP^1,Z)=0$，故任何 $G$-丛都可提升为 $G'$-丛。由于已知 $G'$ 的 (1)，这蕴涵 $G$ 的 (1)。为证 (2)，利用在我们的假设下 $(X_*(Z))^{\pi_1(\PP^1\setminus\{0,\infty\})}$ 无挠，故如 [23] 那样有正合列
\[
0\to X_*(\check Z^{\pi_1(\PP^1\setminus\{0,\infty\})})\to X_*(Z(\ckG')^{\pi_1(\PP^1\setminus\{0,\infty\})})\to X_*(Z(\ckG)^{\pi_1(\PP^1\setminus\{0,\infty\})})\to 0.
\]

为证第三条断言，我们遵循 Ramanathan [33] 或 Faltings [14, Lemma 4] 的论证。

先固定支配权的记号。首先 $X_*(T)^{\pi_1(\mathbb{G}_m)}\otimes\mathbb{Q}=X_*(S)\otimes\mathbb{Q}$。此外，相对根 $\Phi(G,S)$ 张成 $S$ 的在 $G$ 的中心的限制下平凡的特征构成的子空间。记 $X_*(T)^{\pi_1(\mathbb{G}_m),+}\subset\pi_0(\Bun_T)$ 为使对任何正相对根 $a\in\Phi(G,S)^+$ 有 $a(\gamma)\ge 0$ 的那些 $\gamma\in X_*(T)^{\pi_1(\mathbb{G}_m)}$ 构成的子集。

首先证明
\[
G(k((s)))=\bigsqcup_{\lambda\in X_*(T)^{\pi_1(\mathbb{G}_m),+}}G(k[t])\lambda(s)G(k[[s]]).
\]
由 (1) 知每个 $G$-丛在任一点之外平凡。特别地，任何 $G$-丛可约化到 Borel 子群 $B$。设 $E$ 是 $G$-丛，并选其到 $B$ 的约化 $E_B$。

对任何特征 $\alpha:T\to\mathbb{G}_m$，记 $E_B(\alpha)$ 为 $\PP^1$ 上相伴的线丛。由于 $X_*(T)^{\pi_1(\mathbb{G}_m)}\otimes\mathbb{Q}\cong X_*(S)\otimes\mathbb{Q}$，次数 $\deg(E_B(\alpha))\in\mathbb{Q}$ 对 $\alpha\in X_*(S)$ 也有定义。我们断言：若对所有正单根 $a_i\in\Phi(G,S)^+$ 有 $\deg(E_B(a_i))\ge 0$，则丛 $E_B$ 可约化到 $T$。

为证此，记 $U$ 为 $B$ 的幂幺根，$E_T:=E_B/U$ 为诱导的 $T$-丛。为证 $E_B$ 由 $E_T$ 诱导，只需证明 $H^1(\PP^1,E_T\times^TU)=0$。群 $U$ 有滤链，其子商由根子群给出。考虑正相对根 $a\in\Phi(G,S)$。根子群 $U_a$ 是 $[N]^*(\oplus U_{\alpha'})$ 的直和项，其中和跑遍限制到 $S$ 上等于 $a$ 的那些根 $\alpha'\in\Phi(G,T)$。故 $U_a$ 是满足 $H^1(\PP^1,V)=0$ 的向量丛 $V$ 的直和项。

类似地，$E_T\times^TU_a$ 是 $E_T\times^T[N]^*(\oplus U_{\alpha'})$ 的直和项。由于 $[N]^*E_T$ 是 $\pi_1(\mathbb{G}_m)$-不变的，这蕴涵当 $\deg(E_B(a))\ge 0$ 时该丛的 $H^1$ 为 $0$。于是也得 $H^1(\PP^1,E_T\times^TU)=0$。

若约化 $E_B$ 对某个单根 $a_i$ 不满足 $\deg(E_B(a_i))\ge 0$，我们想修改 $E_B$。考虑由 $B$ 与 $U_{-a_i}$ 生成的抛物子群 $P_{a_i}\subset G$。根子群 $U_{\pm a_i}$ 定义了 $P$ 的一个子群 $L$，使 $L$ 的单连通覆盖或同构于 $[n]^*SL_2$，或同构于 $[n]^*SU_3$，其中 $n$ 整除 $N$，见 [7, §4.1.4]。半单商 $P^{ss}/Z(P^{ss})$ 同构于 $L^{\mathrm{ad}}$，且 $P/B\cong L/(L\cap B)$。我们断言结论对 $L^{\mathrm{ad}}$-丛成立，因为这些丛可用向量丛描述。若 $L^{\mathrm{ad}}=[n]^*\mathrm{PGL}_2$，则结论由秩 $2$ 向量丛的相应结论推出。酉群的情形类似，将在下面引理 8.1 中说明。

于是可找到新的到 $B$ 的约化 $E'_B$，使 $E_{P_{a_i}}$ 不变，但 $E'_B(a_i)\ge 0$。这意味着：对基本权 $\epsilon_k$（即由 $B$ 与所有 $j\neq k$ 的 $U_{a_j}$ 生成的极大抛物子群 $P_k$ 的伴随表示的行列式的倍数），当 $k\neq i$ 时 $\deg E_B(\epsilon_k)$ 不变，而 $\deg E'_B(\epsilon_i)$ 大于 $\deg E_B(\epsilon_i)$。

然而 $E'_B(\epsilon_i)$ 是 $\wedge^{\dim(\mathrm{Lie}(P_i))}E(\mathrm{Lie}(G))$ 的子丛，故所有这些线丛的次数有界，于是该过程必然终止于满足对所有单根 $a_i$ 有 $\deg(E_B(a_i))\ge 0$ 的 $B$-约化。这证明了 (3)。

下面推出 Birkhoff 分解 (4)。对常群的情形，这通常由分解 $G=B^-W_0B$ 与 (3) 推出。对一般 $\mathbb G$，其类比由 [7, Theorem 4.6.33] 提供：记 $\mathbb G_0:=\mathbb G^0_{\mathrm{red}}$ 为 $\mathbb G$ 在 $0$ 处纤维的约化商，$B_0\subset\mathbb G_0$ 为 $B$ 的像。记 $U_0$ 为 $B_0$ 的幂幺根，使 $\mathbb G_0=U_0W_0B_0$。所引结果表明 $B_0$ 在 $\mathbb G$ 中的原像正是 $I(0)$。按 $\mathbb G$ 的构造，$U_0(k)$ 的元素可提升到 $U(\PP^1)$。于是
\[
\mathbb G(k[[s]])=U_0W_0I(0).
\]
类似地，按 $\mathbb G$ 的构造，该分解蕴涵 $\mathbb G(\mathbb{A}^1)=I^-(0)W_0U_0$。对支配的 $t\in T(k((s)))$ 与 $b\in N_0$ 有 $t^{-1}bt\in I(0)$。于是利用 (3) 得
\[
\mathbb G(k((s)))=I^-(0)T(k((s)))\mathbb G(k[[s]]).
\]
现以同样方式分解 $\mathbb G(k[[s]])$。对任何 $t$ 可取 $w\in W_0$ 使 $wtw^{-1}$ 支配。再取 $B'_0\subset\mathbb G_0$ 为 $wB_0w^{-1}$，并写 $\mathbb G_0=B'_0W_0B_0$。则可用与前相同的论证推出 (4)（仍假设 $k$ 代数闭）。

下面推出 $k$ 为有限域的情形。先假设 $\mathbb G$ 在全分歧覆盖 $[N]:\PP^1\to\PP^1$ 上分裂。此时嵌入 $\widetilde W\to\mathrm{Gr}_{\mathbb G,x}$ 定义在 $k$ 上，故 $\Bun_{\mathbb G}$ 的所有几何点都定义在 $k$ 上。此外我们断言任何 $\mathbb G(0,0)$-丛的自同构群连通。一旦证明这一点，(1)、(3)、(4) 便由 Lang 定理在 $k$ 上成立。

有 $H^0(\PP^1,T)\hookrightarrow H^0(\PP^1,[N]^*(\mathbb{G}_m))=\mathbb{G}_m$。因此 $H^0(\PP^1,T)\cong T^{0,\mathrm{red}}$，它连通。又对任何 $w\in\widetilde W$，相应的 $k$ 上 $\mathbb G(0,0)$-丛的自同构群是 $I^-(0)\cap wI(0)w^{-1}$。该群以 $H^0(\PP^1,T)$ 为商，其核是仿射根对应的根子群之积，也是连通的。

一般情形由此经 Galois 下降得到，利用 [23, Remark 9]：Iwahori--Weyl 群可计算为 $k$ 的可分闭包上的 Iwahori--Weyl 群的 Galois 不变量。
\end{proof}

上述证明中用到了下述特殊情形。记 $SU_3$ 为覆盖 $[2]:\PP^1\to\PP^1$ 的拟分裂酉群。它可描述为 Hermite 型 $h(x_1,x_2,x_3)=x_1x_3^\sigma+x_2x_2^\sigma+x_3x_1^\sigma$ 的特殊酉群。记 $PSU_3$ 为相应的伴随群。

\begin{lemma}
任何 $PSU_3$-丛 $P$ 都有到 $B$ 的约化 $P_B$，使对正根 $\alpha$ 有 $\deg(P_B(\alpha))\ge 0$。
\end{lemma}

\begin{proof}
定义 $GU_3$ 为由 $SU_3$ 把其中心扩充为 $[2]^*\mathbb{G}_m$ 所得的群，于是有正合列 $1\to[2]^*\mathbb{G}_m\to GU_3\to PSU_3\to1$。同样，每个 $PSU_3$-丛都由某个 $GU_3$-丛诱导。这样的丛可视为覆盖 $\PP^1\xrightarrow{[2]}\PP^1$ 上的秩 $3$ 向量丛 $E$ 带一个取值于形如 $[2]^*L$ 的线丛的 Hermite 型。此时，给出到 $B$ 的约化只需给出各向同性的线子丛 $E_1\to E$，因为它定义旗 $E_1\subset E_1^\perp\subset E$。若 $E$ 非半稳定，则 $E$ 的典范子丛给出正次数的各向同性子丛，故此情形约化存在。若 $E$ 半稳定，则 Hermite 型给出整体同构 $E\xrightarrow{\sim}\sigma^*E^\vee\otimes[2]^*L$。但这样的同构必为常值，故可找到次数为 $\frac{\deg(E)}{3}$ 的各向同性子丛。
\end{proof}
